\scanpage{70}
\begin{center}6 - 32 -\hfill n\textsuperscript{o} 392\end{center}

26) Soient $A$ un anneau factoriel, $K$ son corps des fractions. Pour que des éléments $f_i$, $(1\leq i\leq r)$, de $B=A[X_1,\ldots,X_n]$ soient tels que l'idéal $\sum_{i=1}^r Bf_i$ soit égal à $B$, il faut et il suffit qu'il existe des polynômes $v_i$, $(1\leq i\leq r)$, de $K[X_1,\ldots,X_n]$ tels que $\sum_{i=1}^r v_i f_i=1$ et qui vérifient en outre la condition suivante : si l'on pose $v_i=w_i/d$, où $d\in A$, les polynômes $w_i$ appartiennent à $A[X_1,\ldots,X_n]$ et le p. g. c. d. de l'ensemble des coefficients de tous les $w_i$, $(1\leq i\leq r)$, est égal à 1, alors, pour tout élément extrémal $p$ de $A$ divisant $d$, l'idéal engendré par les classes des $f_i$ dans l'anneau $(A/Ap)[X_1,\ldots,X_n]$ est cet anneau tout entier.

\begin{center}-:-:-:-\end{center}
\margin{69}
